\documentclass[11pt]{article}
\usepackage[margin=1in]{geometry}
\usepackage{amsmath,amssymb,amsthm}
\usepackage[T1]{fontenc}
\usepackage{lmodern}
\usepackage{xcolor}
\usepackage{booktabs}
\usepackage{array}
\usepackage{longtable}
\usepackage[hidelinks]{hyperref}
\usepackage{url}

\newtheorem{theorem}{Theorem}
\newtheorem{lemma}[theorem]{Lemma}
\newtheorem{proposition}[theorem]{Proposition}
\theoremstyle{definition}
\newtheorem{definition}[theorem]{Definition}
\theoremstyle{remark}
\newtheorem{remark}[theorem]{Remark}

% A declaration of our own development. Every \decl{...} is checked by base name against the
% repository by the name gate (Layer 11) before deposit; the gate does not check namespaces.
\DeclareUrlCommand\decl{\urlstyle{tt}\def\UrlBreaks{\do\_\do\.}}
\newcommand{\PH}[1]{\textcolor{red}{\textbf{[PLACEHOLDER:} #1\textbf{]}}}
% File:line anchor into the Paper K repository, public at the pinned commit (see \S\ref{sec:availability});
% the anchors are machine-checked against that commit.
\newcommand{\pkcommit}{464a59df987e149e1e545cbd5299196ecce354ed}
\newcommand{\src}[3]{\texttt{#2}:#3}
\newcommand{\RR}{\mathbb{R}}

\title{Apollonius centres in Lean~4: the cofactor construction, equal Laguerre power, and degenerate configurations}
\author{Josh Bald\\ \small Independent researcher, Ontario, Canada\\ \small\texttt{jpbald93@gmail.com}\\
\small ORCID: \texttt{0009-0002-1317-6489}
\thanks{\copyright~2026 Josh Bald. Licensed under CC BY 4.0
(\url{https://creativecommons.org/licenses/by/4.0/}). The accompanying Lean code is licensed under
the Apache License 2.0.}}
\date{17 November 2026}
\emergencystretch=3em
\begin{document}
\sloppy
\maketitle

\begin{abstract}
We report a Lean~4 formalisation, on top of Mathlib, of part of the Lie-coordinate theory of
Apollonius cycles behind M.~Płatek's concurrency theorem for $n+2$ cycles in $\RR^n$
(arXiv:2604.03067). The formalisation covers three layers. First, the linear algebra: the
Lie form (whose signature $(n+1,2)$ is identified by hand, not in Lean), its null cone, the tangency
criterion, and, at $n=2$, the cofactor
construction that produces from $n+2$ cycles a vector orthogonal to all of them. On that basis we prove the
collinearity step of the source's Proposition~3.1 --- the concurrency point lies on the line joining
the centres of the two Apollonius cycles of each $(n+1)$-subset --- for $n=2$ and for general $n$,
under the hypothesis that the orthogonal complement of the remaining cycles is spanned by the two
Apollonius cycles; the source's derivation of that hypothesis from the existence of exactly two
distinct Apollonius cycles is \emph{not} formalised. Second, two statements the source does not
contain: an explicit determinant formula for the point and radius of equal Laguerre power at $n=2$,
which a hand argument --- not packaged as a theorem of the pinned library, though certified in the supplementary layer-13 implication --- identifies with the concurrency point $P_X$, and the
equal-power property of $P_X$: every finite cycle ($v\neq0$) orthogonal to it has the same Laguerre
power. Third, degenerate configurations where the source's centre-line is
undefined: concentric Apollonius pairs (the square), equilateral circle triples, the regular
tetrahedron in $\RR^3$, and simplices at general $n$ whose centres are equidistant from the origin
(under an equal-norm hypothesis that cannot be dropped: a Lean
counterexample at $n=2$ shows the unqualified statement is false). Two negative results are kept visible: the general-$n$ "centroid" reading of the
degenerate case is false, with a Lean counterexample, and our conjecture that the completion map of the
source's Theorem~3.4 fixes $P_X$, refuted numerically. Theorem~3.4, Theorem~3.6 and the Morita corollary
of the source are \emph{not} formalised in any dimension; we state precisely what would be needed.
The development is gated (73 declarations, the general-$n$ file included, standard axioms only), carries a nine-test tampering
suite and a mutation/ablation/satisfiability battery, and was rebuilt from a clean export on a second
machine with byte-identical sources.
\end{abstract}

\section{Introduction}

Płatek's paper on Apollonius cycles in $\RR^n$ \cite{platek} states a concurrency theorem with an
unusually clean linear-algebra core. Given $n+2$ cycles $X_1,\dots,X_{n+2}$ in $\RR^n$, for each
index $i$ form the two Apollonius cycles $A_i,A_i'$ tangent to all the cycles except $X_i$; then the
$n+2$ lines joining the centres of the pairs $A_i,A_i'$ are concurrent at a point $P_X$
(Proposition~3.1). The point has a second description: it is the centre of the sphere tangent to
every hyperplane tangent to the two cycles of each pair (Theorem~3.6), which in $\RR^3$ specialises
to Morita's inscribed-sphere theorem. The framework is Lie's: cycles are points of a projective
quadric, tangency is orthogonality, and the Apollonius problem becomes linear algebra.

This paper reports what of that programme has been machine-checked in Lean~4 with Mathlib, and what
has not. The division matters more than the count. Our formalisation covers the Lie-coordinate model in $\RR^2$, $\RR^3$ and at general $n$, and the
cofactor construction at $n=2$; it proves the
collinearity step of Proposition~3.1 whenever the relevant orthogonal complement is spanned by the
Apollonius pair; and it proves, at $n=2$, an explicit determinant formula for the equal-power point that the source
does not give. What it does not do is derive the spanning hypothesis from the geometry of the
Apollonius problem, and it does not construct the hyperplane families $T_i$ that Theorem~3.6 is
about. So no theorem of this paper should be read as a formalisation of Theorem~3.4, Theorem~3.6 or
the Morita corollary; \S\ref{sec:not} lists exactly the missing steps.

The paper's own contributions, in the sense that they are not statements of the source, are:
\begin{enumerate}
\item an explicit formula at $n=2$ for the point and radius of equal Laguerre power, as a quotient of
$4\times4$ determinants (Theorem~\ref{thm:explicit}), which a hand derivation identifies with
$P_X$;
\item the equal-power property of $P_X$ (Theorem~\ref{thm:equalpower}), which turns the
concurrency point into a centre of a sphere in the ordinary metric sense and identifies the point in
degenerate configurations where the source's centre-line does not exist;
\item proved degenerate cases: an equilateral triple of equal circles forces $P_X$ to the triangle
centre, independently of any fourth circle (Theorem~\ref{thm:equilateral}); the square configuration has
concentric Apollonius pairs, so the source's line is undefined, while the orthogonal complement is
two-dimensional and every admissible $P$ still projects to the square's centre
(Theorem~\ref{thm:square}); the regular tetrahedron in $\RR^3$ and the regular simplex with equal
centre norms at general $n$ likewise force $P_X$ to the symmetry centre
(\S\ref{sec:degenerate});
\item two negative results, one of them a refutation in Lean of a statement we had conjectured
(Theorem~\ref{thm:negative}), the other a numerical refutation of an invariance property that a
reader of Theorem~3.4 would plausibly expect (\S\ref{sec:negative}).
\end{enumerate}

The verification story is reported as data, not as a claim of correctness: a gate over 73
declarations, a nine-test tampering suite, a mutation/contradiction/ablation/satisfiability battery,
a definition-bridge layer, and 21 claim rows restated in pure Mathlib vocabulary across three
files (\S\ref{sec:verification}). Earlier numerical checks and the general-$n$ proof run exposed two
of our own statements as false --- one of them is now a Lean theorem in its negated form
(\decl{PaperKn.brief_E_false}) --- and the restatement and source-comparison audits exposed the
misleading tetrahedral name.

\paragraph{What is deliberately absent.} No claim of priority: a prior-art scout
(\S\ref{sec:prior}) found no formalisation of any statement of the source, but it did find a
general-$n$ formalisation of the augmented Descartes theorem that pre-empts the \emph{identities} in
the Soddy family; we cite it as a resource and do not claim the identities as new. No claim about
space forms, about the two-step construction (Theorem~3.4), or about hyperplanes as cycles at
infinity: all are outside this paper.

\section{The source and its conventions}\label{sec:source}

We fix Płatek's notation \cite{platek}. A \emph{cycle} in $\RR^n$ is represented by
\[
X=(r,\,c_1,\dots,c_n,\,v,\,w)\in V:=\RR^{n+3},
\]
where $r$ is the signed radius, $c=(c_1,\dots,c_n)$ the centre coordinates, and $[c_1:\cdots:c_n:v]$
is the homogeneous centre, with $v=0$ the points at infinity. Cycles satisfy the relation
\begin{equation}\label{eq:def21}
2wv=r^2-\lVert c\rVert^2 ,
\end{equation}
and the bilinear form of the model is
\begin{equation}\label{eq:sourceform}
(X\mid Y)=-x_1y_1+\sum_{i=2}^{n+1}x_iy_i+x_{n+2}y_{n+3}+x_{n+3}y_{n+2}
=-rr'+c\cdot c'+vw'+wv' .
\end{equation}
A projective point $[X]$ is a cycle iff $(X\mid X)=0$ (Lemma~2.4), and two cycles are tangent iff
$(X\mid Y)=0$ (Lemma~2.5). The form has signature $(n+1,2)$ --- $n+1$ positive directions and two negative, the $r$ slot and the
one direction of the hyperbolic $(v,w)$ plane, whose other direction supplies the $(n+1)$-st positive
one; the quadric is Lie's.

The results we are concerned with are:

\begin{itemize}
\item \textbf{Proposition~3.1.} For $n+2$ cycles $X$, and for each $i$ the two Apollonius cycles
$A_i,A_i'$ of $X\setminus X_i$, the lines through the centres of $A_i,A_i'$ are concurrent at a
point $P_X$.
\item \textbf{Theorem~3.4.} A two-step construction produces, for each $i$, a further pair
$B_i,B_i'$ whose centre-lines pass through the same $P_X$.
\item \textbf{Theorem~3.6.} With $T_i$ the set of hyperplanes tangent to both $A_i$ and $A_i'$, and
$P_X$ not at infinity, there is a sphere centred at $P_X$ tangent to every hyperplane in every $T_i$.
\item \textbf{Corollary (Morita).} For four mutually externally tangent spheres and an outer sphere
tangent to all four, there is a unique inscribed sphere.
\item \textbf{Corollaries.} The circumcentre, incentre and outer-Apollonius-circle theorems, the
Soddy line and the Soddy gasket, as instances.
\end{itemize}

Two features of the source matter for a formalisation and are worth stating at the outset, because
they determine what a machine proof of Proposition~3.1 has to contain. First, the argument has two
independent halves: a linear-algebra half, in which $P$ is a null vector of the $(n+2)\times(n+3)$
system of orthogonality to all $n+2$ cycles, and a geometric half, in which one shows that the
complement of the $n+1$ cycles other than $X_i$ is \emph{two}-dimensional and is spanned by the two
Apollonius cycles. Only the first half is linear algebra in the strict sense; the second uses the Apollonius
geometry. Second, the concurrency point is defined by a single linear condition --- $[P]$ is the null
vector of a matrix with $n+2$ rows and $n+3$ columns, and at full row rank that null vector is the
matrix's cofactor vector --- and this is what makes the explicit formula of \S\ref{sec:explicit}
possible.

We also record a defect in the source's statement of Theorem~3.4: it quantifies over pairs
$B_1,\dots,B_n$ while the index runs to $n+2$. Nothing in this paper depends on it; we note it so
that no reader takes our silence for agreement.

\section{The encoding in Lean}\label{sec:encoding}

\paragraph{Coordinates.} A cycle at $n=2$ is a function $\texttt{Fin 5}\to\RR$ with slots
$(r,a,b,v,w)$, in that order; at $n=3$ a function $\texttt{Fin 6}\to\RR$ with an extra centre slot;
at general $n$ a structure \decl{PaperKn.Cyc} with fields \texttt{r}, \texttt{c : Fin n $\to$ $\RR$},
\texttt{v}, \texttt{w}. All three carry the same form, written out rather than derived from a matrix:
\begin{equation}\label{eq:leanB}
B(X,Y)=rr'-c\cdot c'-vw'-wv' .
\end{equation}
The Lean definition at $n=2$ is \decl{PaperK.B} (\src{PaperK}{LieCore.lean}{16}); the general-$n$
one is \decl{PaperKn.B}. The two are the same polynomial, and the $\RR^3$ version
(\decl{PaperK3.B}) differs only in the number of centre slots.

\paragraph{The sign relative to the source.} Comparing \eqref{eq:leanB} with \eqref{eq:sourceform},
our $B$ is the \emph{negative} of the source's $(X\mid Y)$. This is not a convention we can choose
away: the null cone and orthogonality are unchanged, so every statement of the form ``$B(X,Y)=0$''
transfers verbatim to the source's form, but every statement about a \emph{value} flips sign. We
state the two consequences once and use them throughout:
\begin{itemize}
\item \emph{Zero-set statements transfer.} Null cone, tangency, and ``$P$ is orthogonal to $X_i$''
are the same conditions in both conventions.
\item \emph{Value statements are stated for $B$.} The equal-power constant is
$-B(P,P)/p_v^2$; in the source's convention it reads $+(P\mid P)/p_v^2$, whose sign is the source's
own.
\end{itemize}
The relation \eqref{eq:def21} appears as \decl{PaperK.B_self_eq_zero_iff}: $B(X,X)=0$ iff
$X_0^2-X_1^2-X_2^2=2X_3X_4$, i.e.\ $r^2-\lVert c\rVert^2=2vw$. At general $n$ the same statement is
\decl{PaperKn.B_self_eq_zero_iff} and reads $r^2-\sum c_k^2=2vw$.

\paragraph{The Gram matrix.} At $n=2$ the form is represented by a symmetric matrix $G$ with
$G^2=1$ and $B(X,Y)=X\cdot(GY)$ (\decl{PaperK.G}, \decl{PaperK.B_eq_dot_G},
\decl{PaperK.G_mul_G}, \decl{PaperK.B_G_mulVec}). Two consequences are used later: the
$G$-adjoint $P\mapsto GP$ is an involution, and a vector $P$ orthogonal to the rows of a matrix is
obtained by applying $G$ to a cofactor vector.

\paragraph{The null cone and tangency.} The tangency criterion is proved in the form the briefs
needed, with both cycles on the null cone:
\begin{equation}\label{eq:tangency}
-2\,v\,v'\,B(X,Y)=(rv'-r'v)^2-\lVert cv'-c'v\rVert^2 ,
\end{equation}
which is \decl{PaperK.tangency_identity}. Two remarks belong here, because both are places where a
formalisation can silently go wrong. First, \eqref{eq:tangency} is \emph{false} as a polynomial
identity in the five free coordinates: the difference between its two sides is
$-\bigl((Y_3)^2B(X,X)+(X_3)^2B(Y,Y)\bigr)$ --- the $v$-slots, with a minus sign --- so the
null-cone hypotheses are not decoration. (The $\RR^3$
statement carries them explicitly; see \decl{PaperK3.tangency_identity}.) Second, for
$v=v'\neq0$ the identity reduces to $2B(X,Y)=\lVert c-c'\rVert^2-(r-r')^2$, which is the ordinary
oriented tangency equation, and the iff form for non-zero $v$-slots is
\decl{PaperK.B_eq_zero_iff_tangent}. At general $n$ both forms are
\decl{PaperKn.tangency_identity} and \decl{PaperKn.B_eq_zero_iff_tangent}.

\paragraph{Cofactors.} For a $4\times5$ matrix $M$ whose rows are four cycles at $n=2$, the cofactor
vector is
\[
\texttt{cofactor}(M)_k=(-1)^k\det\bigl(M\text{ with column }k\text{ deleted}\bigr),
\]
\decl{PaperK.cofactor} (\src{PaperK}{LieCore.lean}{75}). Laplace expansion gives
$M\cdot\texttt{cofactor}(M)=0$ (\decl{PaperK.mulVec_cofactor}), and since $G\cdot v$ is the
$B$-adjoint of $v$,
\begin{equation}\label{eq:cofactororth}
B\bigl(G\cdot\texttt{cofactor}(M),\,M_i\bigr)=0 \qquad\text{for every row } i,
\end{equation}
which is \decl{PaperK.B_G_cofactor_row}. So $P:=G\cdot\texttt{cofactor}(M)$ is the concurrency
candidate, and its affine centre is \decl{PaperK.pt} $P\mapsto(P_1/P_3,P_2/P_3)$
(\src{PaperK}{LieCore.lean}{104}). This is the only place where a division occurs, and every
geometric statement about the centre that we prove assumes $P_3\neq0$, while the coordinate
definitions use Lean's total division; in the source's language, $P_X$ is not at infinity.

\paragraph{The affine chart.} Spheres with finite centre are encoded by
$\texttt{circ}(r,a,b)=![r,a,b,1,(r^2-a^2-b^2)/2]$ at $n=2$ (\decl{PaperK2.circ}) and by the
analogous \decl{PaperK3.sph} and \decl{PaperKn.aff}. These satisfy \eqref{eq:def21} with $v=1$, and
the encoding is the source's Definition~2.1 chart. Hyperplanes ($v=0$) exist as raw vectors at
general $n$ but no theorem is stated for hyperplane cycles or uses $v=0$ non-trivially; this is the
formal gap that keeps Theorem~3.6 out of reach (\S\ref{sec:not}).

\section{The collinearity step of Proposition~3.1}\label{sec:prop31}

\begin{theorem}[Collinearity, $n=2$]\label{thm:prop31}
Let $X_1,\dots,X_4,A,A',P$ be vectors in $\RR^5$ --- no null-cone hypothesis on any of them is
used --- and $i$ an index. Suppose
\begin{enumerate}
\item (spanning) every vector $Y\in\RR^5$ orthogonal to all $X_j$ with $j\neq i$ lies in
$\operatorname{span}\{A,A'\}$;
\item (orthogonality) $P$ is orthogonal to all four $X_j$;
\item (finite centres) the $v$-slots of $A,A'$ and $P$ are non-zero.
\end{enumerate}
Then $P\in\operatorname{span}\{A,A'\}$, and
\[
\exists s,t\in\RR,\quad s+t=1,\quad p(P)=s\,p(A)+t\,p(A'),
\]
where $p$ denotes the affine centre; in particular $p(P),p(A),p(A')$ are collinear.
\end{theorem}

This is \decl{PaperK.prop31} (\src{PaperK}{LieCore.lean}{121}). The proof is short and entirely
linear: $P$ lies in the span by hypothesis~(1), the affine combination is computed by clearing the
$v$-slots,
\[
p(P)=\frac{\alpha A_3}{\alpha A_3+\beta A'_3}\,p(A)+\frac{\beta A'_3}{\alpha A_3+\beta A'_3}\,p(A'),
\]
and collinearity follows from the line-map formula. The weights are the content: they say that the
concurrency point is an affine combination of the two Apollonius centres with coefficients that sum
to one, which is stronger than the bare collinearity of the source's statement.

Two variants are proved. \decl{PaperK.prop31_cofactor} specialises $P$ to the cofactor vector
$G\cdot\texttt{cofactor}(M)$ of \eqref{eq:cofactororth}, so that its \emph{orthogonality} is
discharged by the Laplacian identity instead of assumed. The variant still assumes that the
candidate's $v$-slot is nonzero, and it says nothing when every maximal minor vanishes. Its conclusion
is the collinearity alone, with the span
membership and the affine weights coming from \decl{PaperK.prop31} itself, which the variant applies
rather than re-states. At general $n$ the same theorem is
\decl{PaperKn.prop31}, over $n+2$ cycles, with the conclusion giving both the affine-combination
weights and $\operatorname{Collinear}$ in $\RR^n$.

\paragraph{What is assumed, and why it is assumed.} Hypothesis~(1) is exactly the step the source
derives from the existence of two distinct Apollonius cycles: its $Y_1$ is the two-dimensional
orthogonal complement, and $[A_1],[A_1']$ are two distinct points of $Y_1$. We do not have that
derivation in Lean. It requires (i) formalising the cycle condition on $A,A'$, (ii) formalising
``exactly two distinct Apollonius cycles tangent to $n+1$ given cycles'', and (iii) proving that the
orthogonal complement of the $n+1$ cycles is a projective line which meets the quadric in exactly
$A,A'$. Each is a separate development; the vocabulary for the first is present
(\decl{PaperK.B_self_eq_zero_iff} and the affine chart), but no theorem of ours uses it on $A,A'$,
and the other two are not there. So
Theorem~\ref{thm:prop31} is the linear-algebra core of Proposition~3.1 and not the proposition.

\paragraph{The source's own definition of $P_X$.} One point of the source is easy to miss and is
worth isolating, because it is what makes \S\ref{sec:explicit} possible. In the source, $P$ is
defined by a homogeneous system of $n+2$ equations in $n+3$ unknowns --- the conditions
$(P\mid X_i)=0$ --- so the projective point is the null vector of an $(n+2)\times(n+3)$ matrix. At
$n=2$ that is a $4\times5$ matrix: at full row rank its null space is one-dimensional and spanned by
the maximal-minor (cofactor) vector, while at lower rank that vector vanishes and a kernel vector has
to be chosen separately --- the square configuration of \S\ref{sec:degenerate} is the second case. Our
\decl{PaperK.B_G_cofactor_row} is the coordinate form of that observation, and
\decl{PaperK.prop31_cofactor} is Proposition~3.1's linear step with $P$ taken to be that vector.

\section{An explicit formula for \texorpdfstring{$P_X$}{P\_X}}\label{sec:explicit}

The source gives existence, not a formula. At $n=2$ the cofactor construction can be unwound into
determinants. Write the four cycles as the rows $i\mapsto(r_i,a_i,b_i,s_i,1)$ of an auxiliary data
matrix, with $s_i=a_i^2+b_i^2-r_i^2$ --- these rows are not cycle vectors in the
$(r,a,b,v,w)$ encoding, as the caution below explains --- let $D=\det[r\,a\,b\,1]$ be the $4\times4$ determinant whose columns are
the four rows' $(r,a,b,1)$ entries, and for a column vector $u$ let $\det[r\,a\,b\,u]$ denote the
determinant with the fourth column replaced by $u$.

\begin{theorem}[Explicit equal-power point and radius, $n=2$]\label{thm:explicit}
Assume $D\neq0$. Put
\[
p=\Bigl(\frac{-\det[r\,b\,s\,1]}{2D},\ \frac{\det[r\,a\,s\,1]}{2D}\Bigr),
\qquad
\rho=\frac{-\det[a\,b\,s\,1]}{2D},
\qquad s=a^2+b^2-r^2 .
\]
Then for every $i\in\{1,2,3,4\}$,
\[
\lVert p-c_i\rVert^2-(\rho-r_i)^2=\lVert p\rVert^2-\rho^2+\frac{\det[r\,a\,b\,s]}{D},
\]
that is, the point $p$ with radius $\rho$ has the \emph{same} Laguerre power with respect to all four
cycles.
\end{theorem}

This is \decl{PaperK.explicit_power} (\src{PaperK}{LieCore.lean}{231}) together with
\decl{PaperK.explicit_power_indep} (\src{PaperK}{LieCore.lean}{250}), which states the
right-hand side's independence of $i$. Two cautions, both of which a reader should hold us to.

First, the theorem proves the equal-power property of the displayed point; it does \emph{not} prove
that this point equals $p(P)$ for a cofactor vector. The link is a hand derivation, and the obvious
version of it is \emph{wrong}: the rows $(r,a,b,s,1)$ are not cycles in the $(r,a,b,v,w)$ encoding,
so $G\cdot\texttt{cofactor}$ of that matrix is $B$-orthogonal to the rows and not to the circles,
and its $\texttt{pt}$ is not $p$ (an unseeded random instance gives $(1.266,-4.738)$ against
$p=(1.821,-6.819)$; the instance is not reproducible from the file set, but the discrepancy is
structural and the Layer-13 comparison records the naive variant as refuted). The correct vector, recorded in the Layer-15 comparison, is
$P=(q_0,-q_1,-q_2,2q_3,-q_4)$ with $q=\texttt{cofactor}(M)$: it satisfies
$B(P,\texttt{circ}_i)=0$ by \decl{PaperK.row_relation}, its centre is $p$, and its $v$-slot is
$-2D$. The change of basis and the final identity are not a theorem of the pinned library: they were
derived by hand here, and the slotwise cofactor identities \decl{PaperK.cof0}--\decl{PaperK.cof4} and
\decl{PaperK.row_relation} are their ingredients. The implication itself was subsequently certified in
\texttt{L13Compare.s19\_RB} (\texttt{layer13/BlindCompare\_K.lean}, supplementary and ungated), from the
types of those six declarations.

Second, $\rho$ is the radius that equalises the powers; it is \emph{not} proved to be the inscribed
sphere radius of the source's Theorem~3.6, which is a different statement involving tangent
hyperplanes. The number is the natural candidate --- it is the cofactor vector's $r$-slot divided by its
$v$-slot --- but the identification with an inscribed sphere is not formalised.

The determinant formula was checked numerically before it was proved: on the K1 report's four-circle
instance the four powers agree on the common value $-3.2906$, while the version with the $a$ and $b$
columns interchanged (which our own brief had contained) gives four different values; the closed form
agrees with the cofactor vector to $1.3\times10^{-15}$ on 2000 random configurations
(\texttt{numerics/FORMULA.md}). That check is why the formula in the theorem is the one it is.

\section{Equal power and the metric meaning of \texorpdfstring{$P_X$}{P\_X}}\label{sec:equalpower}

The second ingredient of the extension is that the concurrency point is a point of equal power.

\begin{theorem}[Equal power]\label{thm:equalpower}
Let $P\in\RR^5$ be any vector with $P_3\neq0$ and let $Y$ be a normalised cycle --- a null vector
with $Y_3=1$, in the sense of the affine chart of \S\ref{sec:encoding} --- orthogonal to $P$.
Then the Laguerre power of the point/radius pair $(p(P),P_0/P_3)$ with respect to $Y$ equals
\[
-\frac{B(P,P)}{P_3^{\,2}} .
\]
Consequently, if $P$ is orthogonal to a family of normalised cycles $\{Y_i\}_{i\in\iota}$, then the
powers of $(p(P),P_0/P_3)$ with respect to all $Y_i$ coincide, whatever the index set $\iota$.
\end{theorem}

This is \decl{PaperK.power_eq_of_B_eq_zero} (\src{PaperK}{LieCore.lean}{155}) and
\decl{PaperK.equal_power} (\src{PaperK}{LieCore.lean}{164}). The $n=2$ coordinate form is
\decl{PaperK2.power_eq_const} and \decl{PaperK2.equal_power}
(\src{PaperK}{TwoD.lean}{38},\,\src{PaperK}{TwoD.lean}{47}), which give the common value explicitly as
$\lVert p\rVert^2-\rho^2+2P_4/P_3$; the $\RR^3$ form is \decl{PaperK3.equal_power} with constant
$\lVert p\rVert^2-\rho^2+2P_5/P_4$ (\src{PaperK}{ThreeD.lean}{32}); and the general-$n$ form is
\decl{PaperKn.equal_power}. In the source's convention the constant is
$+(P\mid P)/p_v^2$, which is the sign remark of \S\ref{sec:encoding} in action.

The metric reading is immediate: $\lVert p-c\rVert^2-(\rho-r)^2$ is, up to the sign convention, the
power of $p$ with respect to the circle of centre $c$ and radius $\rho-r$ --- Mathlib's
\texttt{Sphere.power} of that sphere, which is the Laguerre power of the two cycles. So the theorem says: \emph{$P_X$ is the centre of a sphere whose power with
respect to every finite cycle ($v\neq0$) orthogonal to $P_X$ is the same number.} In the degenerate cases of
\S\ref{sec:degenerate} this is the statement that survives when the source's centre-line does not
exist, and it is what the proofs there use.

\section{Degenerate configurations}\label{sec:degenerate}

The source's statements are about the line through two (affine) centres. That line can fail to exist
in two ways: the source's Proposition~3.1 assumes ``exactly two distinct Apollonius cycles'', so it
excludes coincidence of the two but not concentricity; the concentric case occurs below, and there
the concurrency point still has a canonical value.

\paragraph{Equal circles: the circumcentre.} For three circles of a common radius $r_0$ with
non-collinear centres, every vector $P$ orthogonal to all three with $P_3\neq0$ has $p(P)$ equal to
the circumcentre of the three centres. This is \decl{PaperK2.PX_circumcentre} (\src{PaperK}{TwoD.lean}{144}), proved by
identifying the equal-power point with the circumcentre and using uniqueness of the latter
(\decl{PaperK2.equal_power_point_unique}, \src{PaperK}{TwoD.lean}{130}). At $r_0=0$ this recovers
the source's circumcentre example, where the three cycles are point cycles at the vertices; the
credit for the geometric statement is the source's.

\begin{theorem}[Equilateral triple]\label{thm:equilateral}
Let $P\in\RR^5$ be any vector with $P_3\neq0$ orthogonal to the three unit circles centred at
$(0,0)$, $(2,0)$ and $(1,\sqrt3)$ (no null-cone hypothesis on $P$ is required). Then
$p(P)=(1,\sqrt3/3)$.
\end{theorem}

This is \decl{PaperK2.equilateral_PX} (\src{PaperK}{TwoD.lean}{177}), with the parametric version
\decl{PaperK2.equilateral_PX_param} and the direct perpendicular-bisector check
\decl{PaperK2.equilateral_centre_unique}. The point of the theorem is the quantification: the
triangle centre is forced by the three equal circles alone. Their Apollonius pair is concentric ---
the two cycles are distinct but share the triangle centre --- so the centre-line of that subset never
exists, and the point is fixed nonetheless. A fourth circle changes the pairs of the other subsets,
but it neither changes this triple's pair nor moves $p(P)$; the Lean proof does not use the fourth
circle's orthogonality hypothesis at all, and the three-circle statement follows by repeating the
first circle. What the theorem does not assert is existence: if every vector orthogonal to the four has
$v$-slot $0$, there is no finite $P$ to speak of, and a fourth unit circle at $(0,1)$ is an example.
Continuity under a perturbation of one input circle is a numerical observation
(\texttt{numerics/EXTENSION.md} \S2), not a theorem.

\begin{theorem}[Square: concentric Apollonius pairs]\label{thm:square}
Let the four cycles be the unit circles centred at $(\pm1,\pm1)$. For each $3$-subset, the affine
cycles tangent to the three are exactly those centred at the origin with radius $1+\sqrt2$ or
$1-\sqrt2$; the two are distinct and concentric, so no centre-line exists. Nevertheless the space of
\emph{vectors} orthogonal to all four is two-dimensional --- exactly
$\operatorname{span}\{(1,0,0,0,1),(0,0,0,1,1/2)\}$ --- and every such vector with $P_3\neq0$ has
$p(P)=(0,0)$.
\end{theorem}

This is the group \decl{PaperK2.square_apollonius}, \decl{PaperK2.square_apollonius_radii_ne},
\decl{PaperK2.square_null_space}, \decl{PaperK2.square_null_span},
\decl{PaperK2.square_null_indep} and \decl{PaperK2.square_PX} in
\src{PaperK}{TwoD.lean}{217}--\src{PaperK}{TwoD.lean}{286}. Neither spanning vector is null --- their
$B$-values are $1$ and $-1$, a direct computation, not a separate declaration. Two things are worth
separating. The
geometric fact is that the Apollonius pair degenerates to two concentric circles: the source's
Proposition~3.1 hypothesis ``exactly two distinct Apollonius cycles'' \emph{holds} among affine
cycles, so the proposition's conclusion is not merely hard to state, it is undefined at this
configuration --- a gap in the source's hypotheses, exhibited in Lean among affine cycles. That no
oriented line ($v=0$) is an Apollonius cycle here is a hand argument, not a theorem of ours. The algebraic fact is that the
rank drops from four to three, so the null space grows from one dimension to two: $P$ is not unique
as a vector, but the \emph{point} is,
and the equal-power proof (\decl{PaperK2.square_PX_via_equal_power}) shows why, by reducing to the
circumcentre of the triple that omits the fourth circle.

\paragraph{The regular tetrahedron in $\RR^3$.} For the four unit spheres centred at
$(1,1,1),(1,-1,-1),(-1,1,-1),(-1,-1,1)$ and an arbitrary fifth sphere, every vector $P\in\RR^6$
orthogonal to all five with $P_4\neq0$ (no null-cone hypothesis is required) has $p(P)=(0,0,0)$ (\decl{PaperK3.tetra_PX_origin},
\src{PaperK}{ThreeD.lean}{91}), with the uniqueness step \decl{PaperK3.tetra_equidistant_unique}
(the equidistant point is the origin, at squared distance $3$). Two Apollonius cycles of the four
spheres are centred at $P_X$, with radii $\sqrt3+1$ and the orientation-reversed $1-\sqrt3$
(\decl{PaperK3.tetra_E1_apollonius}, \src{PaperK}{ThreeD.lean}{142}); the statement proves that
both are $B$-orthogonal to the four spheres, not that no third exists. The tangency is verified in the
B-form (\decl{PaperK3.tetra_tangent_B}), as two real-number identities
(\decl{PaperK3.tetra_tangent_sq}, \decl{PaperK3.tetra_radii_metric}), and against Mathlib's
\emph{unsigned} tangency predicates through \decl{PaperK3.bridge_sph_isExtTangent}. The scaled family
$\texttt{tetraFamR}$ is the same statement for spheres of common radius $s>0$ centred at
$(\pm s,\pm s,\pm s)$ (\decl{PaperK3.tetraR_PX_origin}), with Apollonius radii
$s(\sqrt3+1)$ and $s(1-\sqrt3)$ (\decl{PaperK3.tetraR_tangent_B}); the family scales the centres as
well as the radius, and the four centres are affinely independent
(\decl{PaperK3.tetra_affineIndependent}).

This theorem carried the misleading name \texttt{PaperK3.tetra\_thm36} in an earlier version, which is
how it invited the reading that Theorem~3.6 had been formalised. It has been renamed
\decl{PaperK3.tetra_E1_apollonius}, because that is what it is: a statement about where the
Apollonius pair of the four spheres is centred, with no hyperplane family and no $T_i$ anywhere in
its statement. The rename is a correction of the record, not a change of claim. Two residues are recorded rather
than tidied: the theorem's own docstring still reads ``Thm 3.6 conclusion in the regular-tetrahedron
case'' and is queued for the next commit, and the Layer-15 comparison was run on the pre-rename file,
whose only difference is the name.

\paragraph{The simplex at general $n$.} The pattern is not special to dimension three. Let
$n+1$ unit spheres have affinely independent centres with centroid $0$. If in addition the centres
have equal norms, then every vector $P$ orthogonal to them with $P_v\neq0$ has $p(P)=0$
(\decl{PaperKn.simplex_E1}; the statement carries an extra cycle hypothesis, unused in the proof), and
the common squared distance from $p(P)$ to the centres equals the common squared vertex norm
$\sum_k (c_{i,k})^2$ (\decl{PaperKn.simplex_circumradius}), which is the coordinate form of the
squared circumradius; the identification with Mathlib's \texttt{Affine.Simplex.circumradius} is made
by hand, not in the pinned library. Given equal norms, the centroid hypothesis is redundant --- the
origin is then the circumcentre --- but the Lean statement keeps it, and with $n+1$ centres in
$\RR^n$ equidistant from the origin the simplex need not be regular for $n\ge3$. The equal-norm hypothesis is not
cosmetic: without it the statement is false, and a counterexample is a theorem
(\S\ref{sec:negative}). This is also the right general statement of the pattern, in the light of the
equal-circles theorem above: $P_X$ is equidistant from the centres at general $n$
(\decl{PaperKn.dist_sq_const}, for spheres of any common radius), hence the circumcentre; that
identification is proved as such only at $n=2$, and the centroid is the circumcentre only when the two
coincide.

\section{Negative results}\label{sec:negative}

Two claims that we had reason to believe, and that are false; both are recorded because a paper that
reports only what closed misleads about what the mathematics is like.

\begin{theorem}[The centroid reading is false]\label{thm:negative}
Let $n+1$ unit spheres in $\RR^n$ have affinely independent centres with centroid $0$. It does
\emph{not} follow that $P_X=0$. At $n=2$ the following is a counterexample.
\end{theorem}

\noindent The natural generalisation of the equilateral and tetrahedral cases is the statement just
refuted. It fails at $n=2$. Take centres $(-1,-1),(2,-1),(-1,2)$ (sum zero,
affinely independent), unit radii, and the fourth cycle $X_0$ of radius $3$ centred at $(3,3)$. The
vector $P=(0,(1/2,1/2),1,3/2)$ is orthogonal to all four --- $X_0$ is not a fifth cycle, there are
only four at $n=2$ --- has $P_3=1$, and $p(P)=(1/2,1/2)\neq0$
(\decl{PaperKn.brief_E_false}, \decl{PaperKn.cex_orth}, \decl{PaperKn.cex_pt_ne}). The
correct statement is the circumcentre one; the failure is exactly the failure of
``circumcentre $=$ centroid''. That the hypothesis is the right repair is proved by
\decl{PaperKn.simplex_E1} in the presence of \texttt{hnorm}.

\paragraph{Invariance under the completion map.} Theorem~3.4 derives, from the Apollonius pairs
$A_i,A_i'$, a second family $B_i,B_i'$ whose centre-lines pass through $P_X$. A reader will ask
whether the construction is invariant, i.e.\ whether the concurrency point of the derived
configuration is again $P_X$. It is not. In six generic planar configurations where every step is
well defined, the distance between the two points is $5.7\times10^{-2}$ to $3.1\times10^{-1}$, while
in the same runs the residual of Theorem~3.4's own conclusion is at most $8.7\times10^{-14}$
(\texttt{numerics/iterate.py}; \texttt{numerics/EXTENSION.md} \S1). This is a numerical
refutation, not a Lean theorem, and it is labelled as such: it kills a plausible claim before a
machine spent budget on it.

A third, smaller negative result belongs here too: our own first general-$n$ brief for the
degenerate case omitted the equal-norm hypothesis, and the error was found by the runner, not by us.
That history is why the hypothesis is stated prominently in \S\ref{sec:degenerate} rather than
buried in a proof.

\section{Verification}\label{sec:verification}

The development is a single Lean~4 project on Mathlib, pinned at Lean \texttt{v4.33.1} and Mathlib
\texttt{0df444a3...} (\texttt{lake-manifest.json}). Every figure below is
read from one commit, \texttt{\pkcommit}: the gate and tampering lines were produced there, the
battery was run there, and the second-machine rebuild was made from a \texttt{git archive} of it.
Where a number comes from an earlier run, the sentence that quotes it says so.

\paragraph{The gate.} \texttt{gate.sh} compiles the project and then checks, for each of 73 required
declarations (24 core statements, 36 definition bridges and 13 general-$n$ theorems), that it depends
only on the standard axioms. It prints \texttt{PASS (73 declarations, standard axioms only)}. The
gate checks axioms, not meaning: it says that these 73 pinned declarations have no \texttt{sorry},
no axiom of their own and no native-evaluation shortcut in their dependency cones. The other
theorems of the files get no separate axiom report: the ones a required declaration uses are covered
transitively by its dependency cone, and all are covered by the forbidden-token scan and the build.
The count above should be read that way.

\paragraph{The tampering suite.} Nine negative tests are applied to copies of the development, each
expected to be \emph{rejected}: a \texttt{sorry}, an axiom declaration split across two lines, a fake
evaluation trick, a macro disguise, a wrapper whose proof goes through a non-standard axiom, and four semantic defects
planted in real statements --- dropping the second null hypothesis from the tangency identity
(\texttt{falsehyp}), changing the null-cone constant from $2$ to $3$ (\texttt{falsestmt}), dropping
the non-collinearity hypothesis $D_3\neq0$ from the planar equal-power uniqueness
(\texttt{falsehyp2}), and changing the tetrahedral radius $\sqrt3+1$ to the false $\sqrt3+2$
(\texttt{falsestmt2}). All nine are rejected, and the suite prints
\texttt{ALL TAMPERING TESTS REJECTED}. The four semantic tests are the ones worth having: each is a
statement a reader might plausibly have written, and each fails at the build stage. That rejection is
of the modified statement/proof pair, not a falsity certificate: two of the tests delete a hypothesis
that the retained proof still uses, so the build fails for a syntactic reason as well. The suite's own
comments give one explicit counterexample and explain a second (for instance
$X=(1,0,0,1,\tfrac12)$, $Y=(1,0,0,1,0)$ for the tangency identity); the other two mutations were
checked by direct computation during review, not by Lean.

\paragraph{The battery.} A nine-layer battery (baseline build, gate, tamper, mutations,
contradiction probe, ablation, linter, satisfiability, bridge inventory) is run at \texttt{\pkcommit}
over the five library files, a superset of the 73 the gate pins. Its numbers: the baseline build compiles 5
of 5 files; 444 of 460 mutants are killed, 16 survive; the contradiction probe finds 0 contradictory
statements in 75 probes; satisfiability is certified for 73 of 127 declarations, with 54 unresolved
and none invalid; ablation needs 238 hypotheses and reports 4 unused, one of them
\decl{PaperKn.simplex_E1}'s extra cycle hypothesis (\texttt{\_hP0}, underscore-marked in the source
because the extra cycle plays no part in the proof); the battery's linter layer reports 0 warnings in
the library, while Lean's own style linter reports one \texttt{unnecessarySeqFocus} warning at
\texttt{TwoD.lean:272}, from the Layer-15 recompile of the same file (byte-identical at
\texttt{a77e2e53...}). Eleven of the 16 survivors are mutants of
\decl{PaperK3.tetra_centre_norm} (three) and \decl{PaperK3.tetra_affineIndependent} (eight),
which is a fair description of where our statement-level redundancy is; one is a helper of the
counterexample of \S\ref{sec:negative} whose mutation keeps the helper true while breaking the
theorem that uses it. They are published by name in the battery report rather than rounded into a
percentage. One scope note belongs with the totals: the battery reports 148 theorem declarations and
scores 127 of them for satisfiability, while the five files contain 121 \texttt{theorem} and
\texttt{lemma} declarations on a direct count, and its linter layer also scanned one scratch file
(\texttt{scratch/paperK\_l14b.lean}, nonzero exit). The report does not itemise the extra
declarations, so the mutation, satisfiability and ablation figures are quoted as that tool's counts
over its scan set, not as a per-declaration audit of the gate's 73. Two further results belong beside
the numbers rather than hidden: the bridge inventory
(Layer~9) reports 24 notions, 3 bridged and 21 without a bridge, and its name-pattern matcher does
not credit the \texttt{bridge\_*} lemmas, so that layer is uninformative rather than negative; and no
mutants were generated for three theorems this paper headlines
(\decl{PaperK.explicit_power}, \decl{PaperK.explicit_power_indep},
\decl{PaperK.prop31_cofactor}), so mutation coverage is not uniform across \S\ref{sec:explicit}.

\paragraph{Definition bridges.} \texttt{PaperK/Bridge.lean} (438 lines) proves 36 statements
connecting the development's vocabulary to Mathlib's: $B$ versus \texttt{dist} and
\texttt{Sphere.power}, the cofactor versus \texttt{Matrix.det}, \texttt{circ}/\texttt{sph} versus
\texttt{EuclideanGeometry.Sphere} and its tangent predicates, and determinant collinearity versus
\texttt{Collinear}. Only two bridges are \texttt{rfl}. The bridges are where the honest gaps are
recorded: Mathlib's tangency predicates are unsigned, so oriented (negative) radii can only appear as
the negated radius and the orientation survives only in the $B$-statements; and the
\texttt{ccx}/\texttt{ccy} formula versus \texttt{Affine.Simplex.circumcenter} is bridged on the
distance half only.

\paragraph{Layer 14: restatements in Mathlib vocabulary.} Independent files, importing
\texttt{Mathlib} only and re-proving the claims, are the closest thing to a statement-level audit
that this workflow has. Part~1 covers seven planar and linear-algebra claims plus one $\RR^3$ claim: 6 equivalent, 1 weaker
(\decl{PaperK.prop31} loses the explicit affine coefficients) and 1 expressed against Mathlib's
circumcentre instead of the Cramer formula; it also found that the tetrahedral statement needs
neither the fifth sphere nor $P_4\neq0$. Part~2's eight remaining planar claims now all compile,
after four proof repairs with no statement changed. Part~3 (the five $\RR^3$ claims) gives 5
equivalent restatements, with the caveat that the tetrahedral theorem is a faithful restatement of
an \emph{E1} fact and not of Theorem~3.6. The layer's credit is mixed and worth stating exactly: the
misleading name was caught by part~3's audit and by the Layer-15 comparison, while the equal-norm
hypothesis was caught by the general-$n$ run and the determinant sign by the numerical check of
\S\ref{sec:explicit}, both before this layer existed.

\paragraph{Layer 15: comparison with the source.} A row-by-row comparison against the verbatim
arXiv source (17 rows, 44 identifiers, 16 quoted statements) is in
\texttt{layer15/L15\_COMPARE.md}. Its conclusion is the reason this paper's abstract is
as narrow as it is: \emph{not one row is equivalent}. The formalised statements are either weaker
than the source's (the spanning hypothesis instead of its derivation), restricted to the affine chart
($v=1$), or about different objects (the $B$-value instead of the projective form). The comparison
also lists, explicitly, the source items with \emph{no} counterpart at all: the decomposition lemma,
the reflection $R_P$ and its interchange lemma, Theorem~3.4, Theorem~3.6, and the Morita corollary.

\paragraph{Layer 13: blind re-formalisation.} The layers above test proofs; this one tests whether the
statements say what the prose says. A model in a separate session, instructed not to look at the Lean
development or any layer report, was given the frozen draft as text and asked to work from its
mathematical prose alone --- it stripped the markup itself and recorded the prose it used --- and
wrote 43 statements in Mathlib vocabulary. This is source-isolation by instruction, not a
pre-sanitised input: the draft it read still contained the paper's declaration names, which it was
told to ignore. That file elaborates with \texttt{rc=0} and no proofs.
A second model, different from the writer, then paired each statement with the repository declaration
I designated for it --- the pairing was made by me from the paper's own theorem labels, not by the
comparer --- and proved the two implications in Lean, using only unfolding, encoding transport
($\RR\times\RR$ versus \texttt{Fin 2 $\to$ $\RR$}, the product encoding of a cycle versus the
structure, \texttt{lag} versus \texttt{power}) and explicit arithmetic; no theorem of the
development is used as a proof step, only as a type. Thirty-three of the 43 rows carry at least one
implication certificate (25 in both directions, 8 in one), three rows stopped with a recorded
difference, two were rejected at the pairing stage because the designated declaration does not state
the blind claim, and five are explicitly unpaired (no counterpart exists: a signature statement, two
restatements of the tangency defect, the numerical instance of \S\ref{sec:explicit}, and the
``neither spanning vector is null'' claim). Of the differences, \emph{one} is real: the blind reading
of the cofactor variant expected the full conclusion of Theorem~\ref{thm:prop31}, while
\decl{PaperK.prop31_cofactor} concludes the collinearity only. No implication certificate is
supplied for s25, s36 or s40: those comparisons need bridges or a change of configuration beyond the
permitted transport steps, and the audit establishes no strength ordering for them. The remaining
differences are artefacts of the blind encoding, and several of them are cases where the Lean is
\emph{stronger} than
the prose: the blind reading supplied a null hypothesis on $P$ that the repository does not need
(Theorem~\ref{thm:equilateral} and the tetrahedral statement), restricted a family to one universe,
chose unit radii where the repository has a common radius $s$, and fixed the tetrahedral centres where
\decl{PaperK3.tetraFamR} scales them. Two further gaps the exercise exposed were in the prose rather
than the code, and both are repaired in this draft: ``cycle'' in
Theorem~\ref{thm:equalpower} had to mean \emph{null}, since the identity is false for a general
vector, and Theorem~\ref{thm:prop31} did not introduce $P$ in its list of objects, so it was
universally quantified. The comparison file is \texttt{PaperK/layer13/BlindCompare\_K.lean} in the
\texttt{LEAN-verification} supporting repository, outside the code commit and its gate --- it was run
against \texttt{\pkcommit} from an untracked \texttt{scratch/} work copy --- with
\texttt{sha256} \texttt{87a22493...} and the report in
\texttt{layer13/L13\_COMPARE.md}.

\paragraph{Two-machine receipt.} The repository was exported with \texttt{git archive} at the
pinned commit and rebuilt on a second machine (arm64 macOS, same Lean version, separate Mathlib
cache). The build exited 0, the gate printed the same \texttt{PASS (73 declarations)} line, the
tampering suite rejected all nine, and the five library files are byte-identical to the first
machine's (\texttt{sha256} \texttt{24e80351...} Bridge, \texttt{9e28c5a3...} GeneralN,
\texttt{6a19ca20...} LieCore, \texttt{3b326b04...} ThreeD, \texttt{a77e2e53...} TwoD).

\paragraph{What is gated, and what is not.} The general-$n$ development is part of the repository
(\texttt{PaperK/GeneralN.lean}, 294 lines) and its 13 axiom-checked theorems are in the gate, so the
count of 73 covers them; the file holds 32 theorem and lemma declarations in all, and the other 19
receive no separate gate-generated axiom report --- those used by required declarations are covered
transitively by their dependency checks, and all are token-scanned and built. The layer-14
restatement files remain outside the gate, and no number above includes them; two of their results go
beyond the pinned library and are recorded as such --- the Cramer-to-Mathlib circumcentre bridge
(\texttt{L14b}, \texttt{circumcenter\_eq\_cc}) and an iff classifying the tetrahedral affine
Apollonius pair (\texttt{L14c}, \texttt{tetra\_apollonius}) --- and neither is in the 73. We say so here because a count that quietly includes
ungated files is the commonest way this kind of paper overstates itself.

\section{What is not formalised}\label{sec:not}

For the reader who wants to know how far the machine proof reaches, here is the boundary, in the
source's own terms.

\begin{enumerate}
\item \textbf{Theorem~3.4} (the two-step construction): not formalised. Numerically checked at $n=2$
before any dispatch (residual at most $8.7\times10^{-14}$, \texttt{numerics/EXTENSION.md} \S1), and the
invariance question of \S\ref{sec:negative} was tested in the same runs. To formalise it one needs the construction of the
derived cycles $B_i,B_i'$, which depends on the reflection $R_P$ of the Lie quadric --- itself
absent from our development.
\item \textbf{Theorem~3.6} (the inscribed sphere): not formalised in any dimension. The missing
ingredients are precise. (a) Hyperplanes as cycles at infinity: an oriented plane
$\{x:n\cdot x=h\}$, $\lVert n\rVert=1$, is encoded as the null vector $(1,n_1,\dots,n_n,0,-h)$, whose
tangency equation with an affine cycle is $r-n\cdot c+h=0$; no theorem of ours is stated for a
hyperplane cycle. (b) The families
$T_i=\{H: H\text{ null},\ H_0=1,\ \lVert H_c\rVert=1,\ v=0,\ B(H,A_i)=B(H,A_i')=0\}$ must be
constructed for every $i$, which
needs the Apollonius pairs for all $n+2$ subsets simultaneously. (c) The ``$P_X$ not at infinity''
hypothesis must be carried as $p_v\neq0$ throughout, and the sphere centred at $P_X$ must be shown
tangent to each $H\in T_i$. Our \decl{PaperK2.power_eq_const} and
\decl{PaperK3.equal_power} are the equal-power shadow of this argument --- the sphere exists as a
point of equal power --- but a tangent hyperplane is not a power equality, and we do not claim the
step.
\item \textbf{The Morita corollary}: not formalised. Beyond the above, our $\RR^3$ tetrahedral
configuration is \emph{not} Morita's: its four unit spheres are $2\sqrt2$ apart, so they are not
mutually tangent, and no tangent cone or inscribed sphere appears. Calling it Morita's theorem would
be an error of exactly the kind the layer-15 comparison exists to prevent.
\item \textbf{The reflection $R_P$ and the interchange lemma}: absent, as noted.
\item \textbf{Space forms}: optional in the plan, out of scope here.
\end{enumerate}

\section{Prior art and what is new}\label{sec:prior}

A prior-art scout was run \emph{before} the proof dispatch, as the workflow requires, and is in
\texttt{SCOUT.md}. It was a first-pass search --- arXiv, GitHub code search,
Mathlib and the AFP; zbMATH and MathSciNet were not searched --- and its own header requires a
pre-draft re-run, which was not made. What it found: no formalisation of any statement of the source,
and no Lie-sphere-geometry development in any proof assistant; the closest work is a formalisation of
the augmented Descartes theorem at general $n$ \cite{augdescartes}, which contains the
Descartes/Soddy \emph{identities} and their converses, and which pre-empts any claim that those
identities are new. Accordingly this paper claims no identity as new.

Our contributions relative to the source, in the only sense a formalisation paper can defend, are the
following. The equal-power property identifies $P_X$ with the Laguerre radical centre of the oriented
cycles, and we expect that characterisation to be classical in Laguerre geometry; the scout searched
for prior \emph{formalisations}, not for the mathematical literature, so nothing below is claimed as
new mathematics.
\begin{itemize}
\item The Lie-coordinate model for affine cycles in $\RR^2$, $\RR^3$ and $\RR^n$, and, at $n=2$,
the cofactor construction, including the coordinate-level tangency criterion in the form \eqref{eq:tangency} that
a Lean proof can actually use (with the null hypotheses made explicit --- the identity is false
without them, and the error term is identified).
\item The collinearity theorem \S\ref{sec:prop31} with its hypothesis made explicit, which is a
statement about the source's Proposition~3.1 that a reader can check line by line.
\item The explicit determinant formula of \S\ref{sec:explicit}, which is Cramer's rule for that
linear system. The source gives existence only; a formula is the kind of statement a machine proof
can be held to, and the equal-power form makes it checkable without solving the Apollonius problem
again.
\item The equal-power property of \S\ref{sec:equalpower}, and the degenerate configurations of
\S\ref{sec:degenerate}, which are not in the source at all. The square configuration in particular is
a case where the source's hypotheses hold among affine cycles while its conclusion is not even
statable --- an observation about the source's hypotheses, checked here among affine cycles, with the
affine restriction made explicit in \S\ref{sec:degenerate}.
\item The two negative results of \S\ref{sec:negative}.
\end{itemize}
We do not claim to have formalised Płatek's theorems; \S\ref{sec:not} is the accounting.

\section{Availability and reproducibility}\label{sec:availability}

The Lean development is a single repository, \texttt{paperk}, with five library files
(\texttt{PaperK/LieCore.lean}, \texttt{TwoD.lean}, \texttt{ThreeD.lean}, \texttt{Bridge.lean},
\texttt{GeneralN.lean}; 272, 332, 217, 438 and 294 lines), a \texttt{gate.sh} and a nine-test \texttt{tests/tamper.sh}. The
commit this paper refers to is \texttt{\pkcommit}, and the repository is published at that commit:
\url{https://github.com/jpbald93/paperk_apollonius/tree/464a59df987e149e1e545cbd5299196ecce354ed}.
Every \texttt{File:line} anchor in this paper resolves there, and the tree contains the gate, the
tampering suite and the Layer-10 bridge module. The pinned commit, the gate line, the tampering
line and the byte-identical two-machine hashes in \S\ref{sec:verification} remain the
reproducibility statement, which does not depend on the remote staying up.

The supporting material --- the prior-art scout, the layer reports, the layer-13 blind
re-formalisation, the layer-15 comparison against the verbatim source, the two review rounds and their triage,
the independent numerical checks, the battery reports and the second-machine receipts --- is
published as the \texttt{PaperK} folder of
\url{https://github.com/jpbald93/LEAN-verification}. The battery and gate logs are quoted at the
pinned commit; every battery and gate figure quoted in \S\ref{sec:verification}, including the
mutation totals, is from the run at that commit.


\appendix
\section{Statement-level comparison with the source}\label{app:l15}

The table below condenses the Layer-15 comparison (\texttt{layer15/L15\_COMPARE.md}, 17 rows, 44
identifiers, 16 statements quoted from the verbatim arXiv source). Verdicts: \textbf{E} equivalent,
\textbf{O} ours weaker or otherwise different, \textbf{N} no source counterpart or a different object.
The count is 0 E, 6 O, 11 N. This is the table behind the abstract's narrow claim: our formalisation is
the linear core, not the theorem. Rows~12--13 use the theorem's current name;
\texttt{L15\_COMPARE.md} lists the same statement as \texttt{tetra\_thm36} (\S\ref{sec:degenerate}).

\begin{longtable}{p{0.19\linewidth} p{0.045\linewidth} p{0.66\linewidth}}
\toprule
our statement & verdict & reason \\
\midrule
\endhead
\decl{PaperK.B}, \decl{PaperKn.B} & N & $B$ is the negative of the source's form: same null cone, same
orthogonality, values flip sign. \\
\decl{PaperK2.circ}, \decl{PaperK3.sph}, \decl{PaperKn.aff} & O & only the $v=1$ affine chart of the
source's Definition~2.1; hyperplanes exist as raw vectors but no theorem is stated for hyperplane
cycles. \\
\decl{PaperK.B_self_eq_zero_iff}, \decl{PaperK3.B_self_zero_iff},
\decl{PaperKn.B_self_eq_zero_iff} & O & proves the coordinate relation
$r^2-\lVert c\rVert^2=2vw$; ``represents a cycle'' and the projective class are not formalised. \\
\decl{PaperK.tangency_identity}, \decl{PaperK.B_eq_zero_iff_tangent} & O & the identity holds for
all null vectors; the tangency iff additionally needs both $v$-slots nonzero; tangency is the oriented
equation, and the identity is false without the null hypotheses. \\
\decl{PaperK.prop31}, \decl{PaperK.prop31_cofactor} & O & the collinearity step given the spanning
hypothesis; the source's derivation of that hypothesis from two distinct Apollonius cycles is absent;
centres at infinity excluded. \\
\decl{PaperKn.prop31} & O & same shape at general $n$; concurrency of all $n+2$ lines is an immediate
corollary but is not itself a stated theorem. \\
\decl{PaperK.equal_power}, \decl{PaperK.power_eq_of_B_eq_zero} & N & no source counterpart: the
equal-Laguerre-power reformulation is ours. \\
\decl{PaperK.explicit_power}, \decl{PaperK.explicit_power_indep} & N & no source counterpart; needs
$\det[r\,a\,b\,1]\neq0$; the link to $\texttt{pt}(P)$ is hand-derived (\S\ref{sec:explicit}). \\
\decl{PaperK2.equal_power}, \decl{PaperK2.PX_circumcentre}, \decl{PaperK2.equilateral_PX},
\decl{PaperK2.square_PX} & N & equal power for arbitrary affine circle families; the circumcentre,
equilateral and square instances are E1 in $\RR^2$ with equal radii; the continuity statement is
numerical only. \\
\decl{PaperK2.PX_circumcentre} at $r_0=0$ & O & recovers the source's circumcentre example as an
instance; the credit is the source's. \\
\decl{PaperK2.square_apollonius}, \decl{PaperK2.square_null_space} & N & the square's Apollonius pairs
are concentric and distinct, so the source's hypothesis holds \emph{among affine cycles} while its
centre-line is undefined; the exclusion of $v=0$ solutions is a hand argument. \\
\decl{PaperK3.tetra_E1_apollonius} vs Theorem~3.6 & N & different statement: no $A_i$, no $T_i$, no
hyperplane; it is an E1 fact about the Apollonius pair's centre. \\
\decl{PaperK3.tetra_E1_apollonius}, \decl{PaperK3.tetra_tangent_B} vs Morita & N & not Morita's
configuration: the four unit spheres are $2\sqrt2$ apart, not mutually tangent. \\
\decl{PaperK3.tetra_PX_origin}, \decl{PaperK3.tetraR_PX_origin}, \decl{PaperK3.equal_power} & N &
equal power for arbitrary affine families in $\RR^3$; the origin conclusions are E1 for the regular
tetrahedron (and its scaled version) plus an arbitrary fifth sphere. \\
Theorem~3.6 at general $n$ & N & not formalised: no statement names the hyperplane families $T_i$. \\
\decl{PaperKn.simplex_E1}, \decl{PaperKn.simplex_circumradius} & N & E1 at general $n$ under the
equal-norm hypothesis, which cannot be dropped from the universal statement, as the $n=2$
counterexample \decl{PaperKn.brief_E_false} shows. \\
\decl{PaperKn.brief_E_false} & N & refutes our own brief's centroid version; it refutes nothing in the
source and shows the correct reading is the circumcentre. \\
\bottomrule
\end{longtable}

\noindent Source items with no Lean counterpart at all, listed so that nothing is read into silence: the
decomposition lemma, the reflection $R_P$ and its interchange lemma, Theorem~3.4 (numerics at $n=2$ only),
Theorem~3.6, the Morita corollary, and the \S4 examples (incentre, outer Apollonius, mixtilinear, Soddy
line, gasket). We also record a defect in the source's Theorem~3.4 statement: it lists pairs
$B_1,\dots,B_n$ while the index runs to $n+2$.

\section*{Acknowledgements}
The Lean work was carried out on two machines dedicated to this series; the battery harness, the
tampering suite and the layer reports were developed for the companion paper on bicentric
quadrilaterals and the arbelos. The prior-art scout and the layer-15 comparison are the parts of the
workflow that changed this paper the most.

\begin{thebibliography}{9}
\bibitem{platek} M.~Płatek, \emph{On the Inscribed Sphere and Concurrent Lines through the Centers of
Apollonius Spheres in $\mathbb{R}^n$}, arXiv:2604.03067v1 (submitted 3 April 2026).
\bibitem{augdescartes} \texttt{JD-Jones-ASES}, \texttt{augmented-descartes-lean}, a general-$n$
formalisation of the augmented Descartes theorem (2026); consulted as prior art in
\texttt{SCOUT.md}.
\end{thebibliography}

\end{document}
